% !TeX root = ../../Book2.tex
% 纲要标题：## 附录 B　二次型、对合与 Clifford 代数
% 纲要位置：计划纲要.md，第 1778 行。

\chapter{二次型、对合与 Clifford 代数}
\label{b2:app:B}

\begin{chapterabstract}
二次型既决定几何中的正交关系，也能规定结合代数中的乘法。本附录证明 Clifford 代数的有序基定理，计算低维偶代数；在已经建立的 Witt 群上加入张量积，得到 Witt 环及其初步不变量；随后把对称型、交替型和 Hermitian 型解释为矩阵代数上的伴随对合，区分标量倍与换基对伴随运算的不同影响，并计算一个右四元数模上的 Hermitian 例子。最后通过自旋群的平面计算，说明代数构造怎样产生旋转，以及底域为什么会影响这种构造的像。
\end{chapterabstract}

\begin{learninggoals}
\begin{itemize}
\item 用外代数上的算子证明 Clifford 代数的维数，分清普通张量积与分次张量积。
\item 说明 Witt 乘法的良定义，计算判别式，并使用二次扩张的范数判别二元型是否等距。
\item 区分第一类、第二类对合以及正交、辛、酉三种矩阵模型，明确线性与半线性的底域。
\item 在实数、有理数和有限域的低维例子中，计算自旋作用并检查满射结论的适用范围。
\end{itemize}
\end{learninggoals}

取实二维二次型 \(q(x,y)=x^2+y^2\)，让对应的生成元 \(e_1,e_2\) 在一个结合代数里满足 \(e_1^2=e_2^2=1\) 且 \(e_1e_2=-e_2e_1\)。于是 \((e_1e_2)^2=-1\)，偶次部分已出现复数的乘法。这不是任意安放的关系：它们来自 \(v^2=q(v)\) 对 \(v=e_1,e_2,e_1+e_2\) 的要求。二次型的数值因此能决定一个新代数的乘法。

\ProseParagraphBreak

按二次关系整理出的有序乘积能否构成基，决定这个代数是否仍完整地包含原来的向量空间。二次型还有两种不同的代数表达：忽略双曲分量，得到 Witt 环中的类；把保持型的恒等式写成算子的关系，得到自同态代数上的伴随对合。这些构造记录的信息不同，不能彼此替代。对特征不为二的域 \(k\) 上维数至少为一的有限维非退化二次空间 \((V,q)\)，令 \(b_q\) 为归一化极化型，可以把它们并列写成
\[
\begin{tikzcd}[column sep=2.5em, row sep=1.5em]
& \operatorname{Cl}(V,q) \arrow[r] & \operatorname{Spin}(V,q) \\
\text{二次型 }(V,q) \arrow[ur] \arrow[r] \arrow[dr]
& {[q]}\in W(k) & \\
& \bigl(\operatorname{End}_k(V),\operatorname{ad}_{b_q}\bigr) &
\end{tikzcd}
\]
Witt 类忽略双曲分量，保留稳定等距信息；Clifford 代数把二次取值写入非交换乘法，但只知道代数的同构类时，仍须检查能否恢复原型。伴随对合记录型的相似类：中央标量倍不改变伴随运算，换基则使伴随对合内共轭，因此它不保留全部等距信息。自旋群从 Clifford 偶部分中选取满足正规化与范数条件的单位元，再用共轭产生特殊正交变换。

\ProseParagraphBreak

本附录默认底域特征不为二。\secref{b2:sec:B01}允许型退化，其余各节的二次空间均非退化。Clifford 代数由\cref{b2:thm:clifford-universal}的商构造得到；其基证明用外代数的外乘与收缩算子检验线性独立。\cref{b2:thm:quadratic-diagonalization}用于低维计算，\cref{b2:thm:witt-cancellation}与\cref{b2:thm:witt-group}保证 Witt 类有确定的各向异性代表。\secref{b2:sec:B03}中，中心单代数与 Brauer 群用于第一类对合的存在判据，Skolem--Noether 定理把分裂对合写成矩阵伴随公式，四元数则给出未分裂的具体模型。

\ProseParagraphBreak
这里着重把已知的线性代数与非交换代数连接起来。特征二的二次型不变量、一般域上二次型的完整分类、带对合除代数上的 Hermitian 理论，以及自旋群的代数群结构，都留待进一步学习。

\ProseParagraphBreak
\secref{b2:sec:B04}再利用有序基与最高外次数，证明 Clifford 偶部分中的这些共轭作用落在特殊正交群中；实数、有理数和有限域的平面计算将检验它在各自底域点上的像。

% !TeX root = ../../Book2.tex
% 纲要标题：### B.1 Clifford 代数的基、分次与低维结构
% 纲要位置：计划纲要.md，第 1780 行。

% 纲要要点（供目录核对）：
%
% * Clifford 代数的泛性质、有序基、伴随分次及低维偶代数；分次张量积回用 Koszul 符号，以外乘和收缩构造表示、最高外次数检测独立，区别原代数与伴随分次；正式证明非退化情形的体积元平方及整个代数的中心
% #### B.1.1 泛性质与正交直和
% #### B.1.2 有序单项式与维数
% #### B.1.3 一至三维的计算
%

\section{Clifford 代数的基、分次与低维结构}
\label{b2:sec:B01}

把二次取值写成乘法关系 \(v^2=q(v)1\)，会得到多大的代数？\cref{b2:thm:clifford-universal}已经说明这种构造的泛性质，却还没有证明向量空间确实嵌入其中。本节补足这一点，并计算几个能与四元数直接比较的例子。始终设 \(\operatorname{char}k\ne2\)、\(\dim_kV=n<\infty\)，记 \(b=\beta_q/2\)；这里允许 \(q\) 退化。

\subsection{泛性质与正交直和}
\label{b2:subsec:B01:01}

回顾\cref{b2:def:clifford-algebra}，\(C=\operatorname{Cl}(V,q)\) 是张量代数按关系 \(v\otimes v-q(v)1\) 所取的商。线性映射 \(f:V\to A\) 只要满足 \(f(v)^2=q(v)1_A\)，就唯一延伸为含幺代数同态 \(C\to A\)。因此等距映射 \(V\to W\) 诱导 Clifford 代数同态；等距同构及其逆诱导互逆同构。这个结论不需要选择基。

\ProseParagraphBreak
关系混合了次数二与零，却保留次数的奇偶性。故 \(C=C^0\oplus C^1\) 是按 \(\mathbb Z/2\mathbb Z\) 分次的代数，其中 \(C^0\) 是\term[Clifford-ou-zidaishu]{偶 Clifford 代数}[even Clifford algebra]。关系理想由偶元素生成，因而是奇偶分次理想；这解释了为什么商中的两部分仍是直和，而不只是两类生成元的名称。
\symidx{\(\operatorname{Cl}^0(V,q)\)：偶 Clifford 代数}

\ProseParagraphBreak
若 \(u\perp w\)，极化关系给出 \(uw+wu=2b(u,w)=0\)，所以正交分量的向量在 Clifford 代数中反交换。普通张量积中的 \(u\otimes1\) 与 \(1\otimes w\) 却彼此交换，因而需要调整乘法的符号。设 \(k\) 为域，\(A,B\) 为按 \(\mathbb Z/2\mathbb Z\) 分次的 \(k\) 代数。将 \((a\otimes b)(a'\otimes b')\) 的两侧因子分别相乘，要把中间的 \(b\) 搬过 \(a'\)；两个奇元素相遇时产生负号，其他情形不改变符号。这与\eqref{b2:eq:graded-wedge-sign}以及\cref{b2:prop:exterior-direct-sum}证明中使用的 Koszul 符号相同。于是，在向量空间 \(A\otimes_kB\) 上规定
\[
(a\otimes b)(a'\otimes b')=(-1)^{|b|\cdot|a'|}aa'\otimes bb'
\]
对齐次元素成立，再双线性延拓。所得代数记作 \(A\widehat\otimes B\)，称为\term[fen-ci-zhangliang-ji]{分次张量积}[graded tensor product]。三因子相乘时，两种结合方式的符号指数同为 \(|b|\cdot|a'|+|b|\cdot|a''|+|b'|\cdot|a''|\)，所以它满足结合律。
\symidx{\(A\widehat\otimes B\)：按奇偶符号相乘的分次张量积}

\begin{proposition}\label{b2:prop:B-clifford-orthogonal-sum}
设 \(k\) 为特征不为二的域，\((U,q_U)\)、\((W,q_W)\) 为有限维 \(k\) 二次空间。若 \((V,q)=(U,q_U)\perp(W,q_W)\)，则有自然的奇偶分次 \(k\) 代数同构
\[
\operatorname{Cl}(V,q)\cong
\operatorname{Cl}(U,q_U)\widehat\otimes\operatorname{Cl}(W,q_W).
\]
\end{proposition}
\begin{proof}
把 \(u+w\) 送到 \(u\otimes1+1\otimes w\)。两个交叉项符号相反，平方等于 \(q_U(u)+q_W(w)\)，故由\cref{b2:thm:clifford-universal}得到从左到右的同态。反向把 \(a\otimes b\) 送到两因子在 \(\operatorname{Cl}(V,q)\) 中的像之积。正交关系给出 \(uw=-wu\)；把两个单项式的各个向量依次交换，得到 \(ba'=(-1)^{|b|\cdot|a'|}a'b\)，所以反向映射尊重上述乘法。两个复合都在生成向量上为恒等，进而在整个代数上为恒等。
\end{proof}

\subsection{有序单项式与维数}
\label{b2:subsec:B01:02}

取 \(V\) 的任意有序基 \(e_1,\ldots,e_n\)。极化关系给出
\[
e_i^2=q(e_i),\qquad e_je_i=-e_ie_j+2b(e_i,e_j)\quad(j>i).
\]
利用第二式交换逆序字母，再用第一式消去相邻的重复字母，每个字都化为严格递增字的线性组合。可以按字长、同一字长下的逆序数依次归纳，保证这个过程终止。困难在于证明化简后没有额外的线性关系。

\begin{theorem}[Clifford 代数的基]\label{b2:thm:B-clifford-basis}
设 \(k\) 为特征不为二的域，\((V,q)\) 为以 \(e_1,\ldots,e_n\) 为有序基的 \(n\) 维 \(k\) 二次空间，允许 \(q\) 退化，则下列 \(2^n\) 个元素构成 \(\operatorname{Cl}(V,q)\) 的 \(k\) 线性基：
\[
1,\qquad e_{i_1}\cdots e_{i_r}
\quad(1\le i_1<\cdots<i_r\le n,\ 1\le r\le n).
\]
特别地，自然映射 \(V\to\operatorname{Cl}(V,q)\) 单射；当 \(n>0\) 时，偶、奇部分的维数都为 \(2^{n-1}\)。
\end{theorem}
\begin{proof}
生成性已经证明。为检验独立性，我们在已有外积基的 \(\Lambda V\) 上构造表示，并让外次数检测有序单项式的长度。最自然的尝试是令 \(\varepsilon_v(\omega)=v\wedge\omega\)，但外乘的平方为零，当 \(q(v)\ne0\) 时不能直接实现 \(v^2=q(v)1\)。因此再加入一个降一次数的算子，使外乘与它的交叉项补出所需的标量。定义
\[
\iota_v(w_1\wedge\cdots\wedge w_r)
=\sum_{j=1}^r(-1)^{j-1}b(v,w_j)
 w_1\wedge\cdots\wedge\widehat w_j\wedge\cdots\wedge w_r,
\qquad\iota_v(1)=0.
\]
外乘 \(\varepsilon_v\) 添入一个 \(v\)，而 \(\iota_v\) 用 \(b(v,w_j)\) 取值后删去第 \(j\) 个因子，分别使外次数增加一和减少一。右侧对 \(w_1,\ldots,w_r\) 交替且多线性，因而由\cref{b2:prop:exterior-alternating-universal}定义了外幂上的线性算子。两次删去不同位置的项成对抵消，故 \(\iota_v^2=0\)；外乘也满足 \(\varepsilon_v^2=0\)。直接将首个因子 \(v\) 与其余因子分开，有
\[
\iota_v\varepsilon_v+\varepsilon_v\iota_v=b(v,v)\operatorname{id}=q(v)\operatorname{id}.
\]
所以 \(T_v=\varepsilon_v+\iota_v\) 满足 \(T_v^2=q(v)\operatorname{id}\)，且对 \(v\) 线性。由 \(T_{v+w}=T_v+T_w\) 还得到
\[
\begin{aligned}
T_vT_w+T_wT_v&=T_{v+w}^2-T_v^2-T_w^2\\
&=\bigl(q(v+w)-q(v)-q(w)\bigr)\operatorname{id}
=2b(v,w)\operatorname{id}.
\end{aligned}
\]
这正是完整的 Clifford 极化关系；平方条件已足以由泛性质给出 \(C\to\operatorname{End}_k(\Lambda V)\)。

选择 \(1\in\Lambda^0V\) 作为作用对象，是因为长度为 \(r\) 的乘积只有每次选择外乘，才能把次数从零升到 \(r\)。展开这 \(r\) 个算子的乘积时，若其中选了 \(s\ge1\) 次收缩，其余 \(r-s\) 次选外乘，则该项非零时的外次数为 \(r-2s\le r-2\)。因此最高外次数只记录全取外乘的那一项。把假设存在的非零线性关系写成 \(\sum_Ic_Ie_I=0\)，其中 \(I=\{i_1<\cdots<i_s\}\)、\(e_I=e_{i_1}\cdots e_{i_s}\)，并约定 \(e_\varnothing=1\)。取 \(c_I\ne0\) 的项中最大长度 \(r\)，将这个关系送入刚构造的表示，即让算子 \(\sum_I c_I T_{e_{i_1}}\cdots T_{e_{i_s}}\) 作用于 \(1\)，再投影到 \(\Lambda^rV\)，其中空集项的算子取恒等。较短的项也不能贡献这个次数，因此投影后得到
\[
\sum_{|I|=r}c_I\,e_{i_1}\wedge\cdots\wedge e_{i_r}=0.
\]
\cref{b2:thm:exterior-power-basis}说明这些外积线性无关，故所有最大长度的系数为零，与 \(r\) 的选取矛盾。这证明独立性。偶、奇子集数相等，是因为固定一个指标并改变它是否属于子集，给出两类之间的双射。
\end{proof}

以后把 \(V\) 与其在 \(C\) 中的像等同。令 \(F_rC\) 为长度不超过 \(r\) 的向量乘积生成的子空间，并置 \(F_{-1}C=0\)，这就给出按字长的滤过。在伴随分次代数 \(\operatorname{gr}C=\bigoplus_{r\ge0}F_rC/F_{r-1}C\) 中，向量 \(v\) 的次数一像满足平方为零，因为 \(v^2=q(v)1\in F_0C\)。由外代数的商构造，得到自然的分次代数满射
\[
\Lambda V\longrightarrow\operatorname{gr}C,
\qquad
v_1\wedge\cdots\wedge v_r\longmapsto
v_1\cdots v_r+F_{r-1}C.
\]
它是满射，因为 \(\operatorname{gr}C\) 由次数一部分生成。化简字时长度不会增加，故\cref{b2:thm:B-clifford-basis}说明 \(F_rC/F_{r-1}C\) 以长度恰为 \(r\) 的有序单项式的类为基。上述映射把 \(\Lambda^rV\) 的外积基逐一送到这组基，所以每个次数上都为同构。由此得到
\begin{equation}\label{b2:eq:B-clifford-associated-graded}
\operatorname{gr}\operatorname{Cl}(V,q)\cong\Lambda V.
\end{equation}
这是一种 PBW 型现象：关系中的低次项在伴随分次中消失，而有序基保证最高次关系之外没有额外的关系。\cref{b2:prop:pbw-filtered-criterion}中由有序单项式基识别伴随分次的做法与此相同；那里得到多项式代数，这里得到外代数。\eqref{b2:eq:B-clifford-associated-graded}是自然的分次代数同构，并不意味着 \(C\) 本身与 \(\Lambda V\) 同构；例如非零的 \(q(v)\) 使 \(v^2\ne0\)。

\subsection{一至三维的计算}
\label{b2:subsec:B01:03}

一维型 \(\langle a\rangle\) 给出 \(k[t]/(t^2-a)\)。当 \(a=0\) 时这是含非零平方零元的代数；当 \(a\ne0\) 是平方时，它同构于 \(k\times k\)；当 \(a\) 非平方时，它是二次域。单是空间维数相同，不能确定 Clifford 代数的单性。

\begin{proposition}\label{b2:prop:B-low-clifford}
设 \(k\) 为特征不为二的域，\(a,b,c\in k^\times\)，\((u,v)_k\) 表示由 \(i^2=u\)、\(j^2=v\)、\(ij=-ji\) 定义的四元数 \(k\) 代数，则
\[
\operatorname{Cl}\bigl(\langle a,b\rangle\bigr)\cong(a,b)_k,
\qquad
\operatorname{Cl}^0\bigl(\langle a,b\rangle\bigr)
\cong k[t]/(t^2+ab),
\]
并且
\[
\operatorname{Cl}^0\bigl(\langle a,b,c\rangle\bigr)
\cong(-ab,-ac)_k.
\]
\end{proposition}
\begin{proof}
对角基满足 \(e_1^2=a,e_2^2=b,e_2e_1=-e_1e_2\)，这正是\cref{b2:def:quaternion-algebra}的关系；两边均以 \(1,e_1,e_2,e_1e_2\) 为基，所以关系给出的满同态是同构。偶部分以 \(1,e_1e_2\) 为基，且 \((e_1e_2)^2=-ab\)，得到第二式。

三维情形令 \(u=e_1e_2,v=e_1e_3\)。计算得 \(u^2=-ab,v^2=-ac,uv=-vu\)。又 \(uv=-a e_2e_3\)，其中 \(a\ne0\)，所以 \(1,u,v,uv\) 生成全部偶部分。这四个元素由\cref{b2:thm:B-clifford-basis}线性无关，给出第三式。
\end{proof}

例如正定实三维型的偶 Clifford 代数为 \((-1,-1)_{\mathbb R}=\mathbb H\)，正定实二维型的偶部分则是 \(\mathbb C\)。两者都在偶部分中出现，而不是把原二次空间直接当成四元数或复数的乘法代数。

\ProseParagraphBreak
Clifford 代数及其与二次型的联系，可进一步参阅\cite[Chapter ``Clifford algebras'']{Lam2005QuadraticForms}。

\begin{proposition}\label{b2:prop:B-clifford-center}
设 \(k\) 为特征不为二的域，\((V,q)\) 为 \(n\ge1\) 维非退化 \(k\) 二次空间，\(e_1,\ldots,e_n\) 为正交基，\(a_i=q(e_i)\in k^\times\)。记 \(C=\operatorname{Cl}(V,q)=C^0\oplus C^1\)，并令 \(\omega=e_1\cdots e_n\)。则
\[
\omega^2=(-1)^{n(n-1)/2}\prod_{i=1}^n a_i\cdot1,
\]
\[
Z(C)=
\begin{cases}
k\cdot1,&n\;\text{为偶数},\\
k\cdot1\oplus k\cdot\omega,&n\;\text{为奇数}.
\end{cases}
\]
这里后一式为中心的向量空间直和分解。两种情形都满足 \(Z(C)\cap C^0=k\cdot1\)。
\end{proposition}
\begin{proof}
正交关系给出 \(e_i^2=a_i\) 及 \(e_ie_j=-e_je_i\)（\(i\ne j\)）。在 \((e_1\cdots e_n)^2\) 中把第二组基向量依次移到相同指标旁，共产生 \(n(n-1)/2\) 个负号，再将每个 \(e_i^2\) 换成 \(a_i\)，便得到平方公式。

由\cref{b2:thm:B-clifford-basis}，每个元素唯一写成 \(c=\sum_I c_Ie_I\)，其中 \(I\subseteq\{1,\ldots,n\}\)，\(e_I\) 为相应的递增乘积，\(e_\varnothing=1\)。设 \(|I|=r\)。把 \(e_j\) 移过其中的基向量，得到
\[
e_je_I=
\begin{cases}
(-1)^r e_Ie_j,&j\notin I,\\
(-1)^{r-1}e_Ie_j,&j\in I.
\end{cases}
\]
对固定的 \(j\)，\(e_Ie_j\) 是 \(e_{I\triangle\{j\}}\) 的非零标量倍：若 \(j\notin I\)，只需整理次序；若 \(j\in I\)，还要乘以非零的 \(a_j\)。映射 \(I\mapsto I\triangle\{j\}\) 是所有指标子集上的双射，所以在 \(e_jc-ce_j\) 中不同 \(I\) 的系数不能相互抵消。由于特征不为二，\(c_I\ne0\) 的项与 \(e_j\) 交换，当且仅当相应的符号为 \(1\)。

若 \(I\) 非空而且不等于全部指标，可分别取 \(j\in I\) 和 \(j\notin I\)，交换条件将同时要求 \(r-1\) 与 \(r\) 为偶数，这是不可能的。空集始终满足交换条件。全部指标只要求 \(n-1\) 为偶数，恰在 \(n\) 为奇数时成立。\(C\) 由所有 \(e_j\) 生成，因此与它们逐一交换就等价于属于中心，得到所述中心分解。奇数维时 \(\omega\in C^1\)，并与 \(1\) 线性无关；偶数维时中心已是标量。因此中心与偶部分的交在两种情形下都是 \(k\cdot1\)。
\end{proof}

% !TeX root = ../../Book2.tex
% 纲要标题：### B.2 Witt 环、判别式与低维分类
% 本轮补充的纲要细目：
% * 比较同维同判别式而偶 Clifford 代数不同的具体型，并说明偶 Clifford 代数本身也未必决定等距类
% 纲要位置：工作记录/计划纲要.md，第 1607 行。

% 纲要要点（供目录核对）：
% * Witt 环、判别式及低维分类
% #### B.2.1 张量积与 Witt 环
% #### B.2.2 判别式与基本理想
% #### B.2.3 二元型及实、复情形

\section{Witt 环、判别式与低维分类}
\label{b2:sec:B02}

正交直和使二次型的稳定等距类组成 Witt 群。若再让两个型作张量积，就能在这些类上定义乘法；关键是说明被忽略的双曲部分在相乘后仍可忽略。本节仍设 \(\operatorname{char}k\ne2\)，且所有二次空间均有限维、非退化。零维空间也包括在内，其型表示加法零元。

\subsection{张量积与 Witt 环}
\label{b2:subsec:B02:01}

设 \((V,q)\)、\((W,r)\) 的归一化极化型分别为 \(b,c\)。在 \(V\otimes W\) 上定义对称双线性型
\[
(b\otimes c)(v\otimes w,v'\otimes w')=b(v,v')c(w,w'),
\]
并双线性延拓；相应二次型记为 \(q\otimes r\)。这同时说明了非纯张量上的取值；只写 \((q\otimes r)(v\otimes w)=q(v)r(w)\) 还不是对所有向量的完整定义。对角化以后，若 \(q=\langle a_1,\ldots,a_n\rangle\)、\(r=\langle d_1,\ldots,d_m\rangle\)，则张量积的对角系数是全部 \(a_id_j\)，均非零，故它仍非退化。

\begin{lemma}\label{b2:lem:B-hyperbolic-tensor}
设 \(k\) 为特征不为二的域，\(H\) 为 \(k\) 上的双曲平面。对每个 \(n\) 维非退化 \(k\) 二次型 \(q\)，有 \(H\otimes q\cong H^{\perp n}\)。
\end{lemma}
\begin{proof}
由\cref{b2:thm:quadratic-diagonalization}写 \(q=\langle a_1,\ldots,a_n\rangle\)。张量积对正交直和分配，因此只须计算 \(H\otimes\langle a_i\rangle\)。若 \(e,f\) 是满足 \(b_H(e,f)=1\) 的双曲基，\(v_i\) 的二次取值为 \(a_i\)，则 \(a_i^{-1}e\otimes v_i,f\otimes v_i\) 的二次取值均为零，互相配对为一，正是一个双曲基。把这些两维块直和即得结论。
\end{proof}

\begin{theorem}\label{b2:thm:B-witt-ring}
设 \(k\) 为特征不为二的域，\(W(k)\) 为有限维非退化 \(k\) 二次型的 Witt 群。在 \(W(k)\) 上规定
\[
[q]\,[r]=[q\otimes r].
\]
这给出交换含幺环结构，乘法单位为 \(\bigl[\langle1\rangle\bigr]\)。
\end{theorem}
\begin{proof}
若 \(q\perp H^{\perp s}\cong q'\perp H^{\perp t}\)，与 \(r\) 张量后，\cref{b2:lem:B-hyperbolic-tensor}将两边新增部分分别化为 \(s\dim r\) 与 \(t\dim r\) 个双曲平面，故 \(q\otimes r\) 与 \(q'\otimes r\) Witt 等价。交换两因子也能改变 \(r\) 的代表元，所以乘法良定义。

结合、交换与分配同构分别由张量积的重新加括号、交换纯张量及分配映射给出。这些映射都保持上述配对，可在纯张量上直接检查，再由双线性推广。\(k\otimes V\to V\)、\(a\otimes v\mapsto av\) 把 \(\langle1\rangle\otimes q\) 等距地送到 \(q\)。加法及负元已由\cref{b2:thm:witt-group}建立，其中 \(-[q]=[-q]\)，于是各项环公理成立。
\end{proof}

这个环称为\term[Witt-huan]{Witt 环}[Witt ring]。\cref{b2:thm:witt-cancellation}保证它仍保存确定的各向异性型：两个各向异性代表若 Witt 等价，便由\cref{b2:thm:witt-decomposition}的唯一性等距。环的加法不是把所有同维型都认作相同。

\subsection{判别式与基本理想}
\label{b2:subsec:B02:02}

用 \(k^\times/k^{\times2}\) 表示非零元素的平方类群。若 \(B\) 是 \(n\) 维型 \(q\) 的对称矩阵，换基把它变为 \(P^{\mathsf T}BP\)，行列式乘以 \((\det P)^2\)。因此可以定义
\begin{definition}\label{b2:def:B-signed-discriminant}
设 \(k\) 为特征不为二的域，\(q\) 为 \(n\) 维非退化 \(k\) 二次型，\(B\) 为 \(q\) 在任意基下的对称矩阵。由公式
\[
\delta(q)=(-1)^{n(n-1)/2}\det B\pmod{k^{\times2}}
\]
得到的平方类称为 \(q\) 的\term[erci-xing-panbieshi]{判别式}[discriminant of a quadratic form]。这里使用带符号的约定；零维型的行列式规定为一，故 \(\delta(0)=1\)。
\end{definition}
\symidx{\(\delta(q)\)：二次型的带符号判别式}

某些文献把不带符号的 \(\det B\) 称为判别式。下面所有公式均采用上述约定，特别是双曲平面满足 \(\delta(H)=1\)。

\begin{proposition}\label{b2:prop:B-discriminant-witt}
设 \(k\) 为特征不为二的域，\(q,r\) 为有限维非退化 \(k\) 二次型，\(\dim q=n\)、\(\dim r=m\)，则上述带符号判别式 \(\delta\) 在 Witt 等价下不变，且
\begin{equation}\label{b2:eq:B-discriminant-sum}
\delta(q\perp r)=(-1)^{nm}\delta(q)\delta(r).
\end{equation}
维数模二定义环同态 \(W(k)\to\mathbb Z/2\mathbb Z\)。其核 \(I(k)\) 是一个理想，且判别式限制为群同态
\[
\delta:I(k)\longrightarrow k^\times/k^{\times2},
\qquad I(k)^2\subseteq\ker\delta.
\]
\end{proposition}
\begin{proof}
正交直和的矩阵分块对角，行列式相乘。指数差
\[
\frac{(n+m)(n+m-1)}2-\frac{n(n-1)}2-\frac{m(m-1)}2=nm
\]
给出\eqref{b2:eq:B-discriminant-sum}。取 \(r=H\)，因 \(m=2\) 且 \(\delta(H)=1\)，得到加一个双曲平面不改变判别式；由 Witt 等价的定义，判别式于是对类良定义。

加双曲平面不改变维数奇偶，而直和使维数相加、张量积使维数相乘，所以维数模二是环同态。在其核中，\eqref{b2:eq:B-discriminant-sum}的符号为一，得到所述群同态。最后令 \(n,m\) 均为偶数。张量积的对角系数给出
\[
\det(q\otimes r)=(\det q)^m(\det r)^n,
\]
它是平方；\(nm\) 被四整除，又使 \((-1)^{nm(nm-1)/2}=1\)。所以 \(\delta(q\otimes r)=1\)。\(I(k)^2\) 中元素是有限个这类乘积的和，利用刚证的群同态性，判别式仍为一。
\end{proof}

\(I(k)\) 称为 Witt 环的\term[Witt-jiben-lixiang]{基本理想}[fundamental ideal]。维数奇偶和判别式提供前两层容易计算的信息，但判别式在整个 \(W(k)\) 上一般不是加法群同态，式中的 \((-1)^{nm}\) 不能删掉。例如 \(\delta\bigl(\langle1\rangle\bigr)=1\)，而 \(\delta\bigl(\langle1,1\rangle\bigr)=-1\)。
\symidx{\(I(k)\)：Witt 环的基本理想}

\subsection{二元型及实、复情形}
\label{b2:subsec:B02:03}

一维非退化型由系数的平方类分类。二元型 \(\langle a,b\rangle\) 则有 \(\delta=-ab\)。若这个平方类为一，取 \(c^2=-ab\)，则 \((c/a,1)\) 是非零各向同性向量，\cref{b2:lem:hyperbolic-plane-split}说明整个二元型就是双曲平面。反过来双曲平面的判别式为一。因而二元情形的判别式已能判断各向同性，却未必能区分所有各向异性型。

\begin{proposition}\label{b2:prop:B-binary-norm-classification}
设 \(k\) 为特征不为二的域，固定非平方 \(d\in k^\times\)，令 \(K=k(\sqrt d)\)，并记
\[
q_a=\langle a,-ad\rangle\qquad(a\in k^\times).
\]
判别式平方类为 \(d\) 的每个二元非退化 \(k\) 二次型都等距于某个 \(q_a\)。对任意 \(a,a'\in k^\times\)，还有
\[
q_a\cong q_{a'}\quad\Longleftrightarrow\quad
a'/a\in N_{K/k}(K^\times),
\qquad N_{K/k}(x+y\sqrt d)=x^2-dy^2.
\]
\end{proposition}
\begin{proof}
设型为 \(\langle a,b\rangle\)，且 \(-ab=ds^2\)。把第二个基向量乘以 \(a/s\)，其系数变为 \(-ad\)，得到所需标准式。

把 \(q_a\) 看成 \(K\) 上的二次取值 \(z\mapsto a\cdot N_{K/k}(z)\)。若 \(a'/a=N(c)\)，乘 \(c\) 是 \(K\) 的 \(k\) 线性双射，且 \(a\cdot N(cz)=a'\cdot N(z)\)，所以给出 \(q_{a'}\to q_a\) 的等距。反过来，等距映射把 \(q_{a'}\) 中取值为 \(a'\) 的向量 \(1\) 送到某个非零 \(z\in K\)，于是 \(a'=a\cdot N(z)\)。范数公式来自 \((x+y\sqrt d)(x-y\sqrt d)=x^2-dy^2\)。
\end{proof}

\begin{example}\label{b2:exm:B-rational-norm-obstruction}
\(\mathbb Q\) 上的 \(\langle1,1\rangle\) 与 \(\langle3,3\rangle\) 判别式相同，却不等距；方程 \(x^2+y^2=3\) 没有有理数解。
\end{example}
\begin{proof}
由\cref{b2:prop:B-binary-norm-classification}，若两型等距，则三是 \(\mathbb Q(i)\) 的范数，即存在有理数 \(x,y\) 使 \(x^2+y^2=3\)。清去分母并约去公因子，得到互素整数 \(X,Y,Z\)、\(Z\ne0\)，满足 \(X^2+Y^2=3Z^2\)。模三可知 \(X,Y\) 都被三整除，代回后 \(Z\) 也被三整除，矛盾。
\end{proof}

在实数域上，\cref{b2:prop:witt-real-complex}已将 Witt 群用符号差识别为 \(\mathbb Z\)。张量积中同号项给正项、异号项给负项，因此符号差相乘，得到环同构 \(W(\mathbb R)\cong\mathbb Z\)，其中 \(I(\mathbb R)=2\mathbb Z\)。复数域的每个非退化型只由维数决定，维数奇偶同时尊重加法和乘法，故 \(W(\mathbb C)\cong\mathbb Z/2\mathbb Z\)。这些简洁分类依赖底域；一般域上的高维分类不能从判别式单独推出。

\begin{example}\label{b2:exm:B-discriminant-even-clifford}
实三维二次型 \(q=\langle1,1,1\rangle\) 与 \(r=\langle1,-1,-1\rangle\) 有相同的判别式平方类，却有不同的偶 Clifford 代数：
\[
\delta(q)=\delta(r)=-1\pmod{\mathbb R^{\times2}},\qquad
\operatorname{Cl}^0(q)\cong\mathbb H,\quad
\operatorname{Cl}^0(r)\cong M_2(\mathbb R).
\]
\end{example}
\begin{proof}
两个对角矩阵的行列式均为一，维数三所给的符号为 \((-1)^3=-1\)，故判别式相同。\cref{b2:prop:B-low-clifford}给出三维公式
\[
\operatorname{Cl}^0\bigl(\langle a,b,c\rangle\bigr)
\cong(-ab,-ac)_{\mathbb R}.
\]
代入便分别得到 \((-1,-1)_{\mathbb R}=\mathbb H\) 与 \((1,1)_{\mathbb R}\)。后一代数的两个生成元可取为 \(\operatorname{diag}(1,-1)\) 和 \(\begin{psmallmatrix}0&1\\1&0\end{psmallmatrix}\)；它们满足平方为一且反交换，并连同单位、二者之积构成 \(M_2(\mathbb R)\) 的基。因此后一代数分裂，前一代数则由\cref{b2:prop:quaternion-division-or-split}及正定范数为除代数，二者不同构。
\end{proof}

等距映射由 Clifford 代数的泛性质延伸为代数同构，并保持奇偶次数；其逆同样延伸，故偶部分的同构类确是型的等距不变量。上例说明它有时能补充判别式的信息。反过来，相同的偶代数也未必给出相同的型，见\cref{b2:ex:B-clifford-distinguishes-ternary}；不能把一个有效的区别方法直接当作完整分类。

% !TeX root = ../../Book2.tex
% 纲要标题：### B.3 带对合的中心单代数
% 本轮补充的纲要细目：
% * 证明相似型对应共轭的伴随对合，解释中央标量倍为何不改变伴随运算
% * 在右四元数模的明确约定下计算 Hermitian 型及伴随对合，保留标量的次序
% * 用 Goldman 元与极大右理想证明第一类对合的存在性判据
% 纲要位置：工作记录/计划纲要.md，第 1617 行。

% 纲要要点（供目录核对）：
% * 带对合的中心单代数
% #### B.3.1 第一类与第二类对合
% #### B.3.2 矩阵代数上的伴随对合
% #### B.3.3 四元数的共轭与正交对合

\section{带对合的中心单代数}
\label{b2:sec:B03}

矩阵转置、共轭转置及四元数共轭都反转乘法次序，但它们对中心的作用不同。先分清这一点，才能把双线性型与 Hermitian 型放进同一个代数问题。对合的定义适用于任意特征；本节后续的具体型与分类计算要求中心域的特征不为二。代数仍有限维、结合且含幺。

\subsection{第一类与第二类对合}
\label{b2:subsec:B03:01}

对于一个非退化的对称型、交替型或 Hermitian 型，伴随算子把等式 \(b(Xv,w)=b(v,X^\dagger w)\) 中的线性作用从第一个变量移到第二个变量。连续移动两个算子时，次序变为相反次序：\((XY)^\dagger=Y^\dagger X^\dagger\)。若希望在未分裂的中心单代数中研究同样的运算，就不能依赖某一组矩阵坐标，而应直接规定这种反乘法性质及取两次恢复原元素的条件。

\begin{definition}\label{b2:def:B-algebra-involution}
设 \(K\) 为域，\(A\) 为有限维中心单 \(K\) 代数。如果可加映射 \(\sigma:A\to A\) 满足
\[
\sigma(1)=1,\qquad \sigma(xy)=\sigma(y)\sigma(x),\qquad
\sigma^2=\operatorname{id},
\]
就称 \(\sigma\) 为 \(A\) 上的\term[daishu-duihe]{对合}[involution on an algebra]。若 \(\sigma|_K=\operatorname{id}_K\)，就称其为第一类对合；否则称为\term[dierlei-duihe]{第二类对合}[involution of the second kind]。
\end{definition}

中心被对合保持：若 \(z\) 与全部元素交换，则 \(\sigma(z)\) 也与全部 \(\sigma(x)\) 交换，而 \(\sigma\) 满射。因此 \(\tau=\sigma|_K\) 是满足 \(\tau^2=1\) 的域自同构，且
\[
\sigma(ax)=\tau(a)\sigma(x)\qquad(a\in K).
\]
第一类对合是 \(K\) 线性的。第二类对合只在固定域 \(k=K^\tau\) 上线性，在 \(K\) 上是半线性的；例如复矩阵的共轭转置是实线性的，通常不是复线性的。

\ProseParagraphBreak
在 \(M_2(\mathbb C)\) 上，转置与共轭转置给出两种具体计算。对 \(X=iE_{12}\)，有
\[
X^{\mathsf T}=iE_{21},\qquad
\overline X^{\mathsf T}=-iE_{21}.
\]
两种运算取两次都恢复 \(X\)，并都反转乘积次序；但转置固定中心元素 \(iI_2\)，共轭转置把它送到 \(-iI_2\)。所以前者属于第一类，后者属于第二类，其固定中心域为 \(\mathbb R\)。区分类别所依据的是对中心的作用，而非矩阵是否含有非实条目。

\begin{proposition}\label{b2:prop:B-involution-center}
设 \(K\) 为特征不为二的域，\(A\) 为有限维中心单 \(K\) 代数，\(\sigma\) 为 \(A\) 上的对合。若 \(\sigma\) 为第二类，记 \(k=K^{\sigma|_K}\)，则 \([K:k]=2\)。若 \(\sigma\) 为第一类，则它给出 \(K\) 代数同构 \(A\cong A^{\mathrm{op}}\)，从而 \(2[A]=0\) 于 \(\operatorname{Br}(K)\)。
\end{proposition}
\begin{proof}
第二类时取 \(t\in K\) 使 \(\tau(t)\ne t\)，令 \(u=t-\tau(t)\)，则 \(u\ne0\)、\(\tau(u)=-u\)。对每个 \(x\in K\)，有唯一分解
\[
x=\frac{x+\tau(x)}2+
u\frac{x-\tau(x)}{2u}.
\]
右侧两个系数均被 \(\tau\) 固定，所以 \(K=k\oplus ku\)。其中 \(u\notin k\)，因为 \(2u\ne0\)。故扩张次数为二。

第一类时，对合反转乘法，恰好是到反代数的 \(K\) 代数同构。\cref{b2:thm:brauer-group}说明反代数代表 Brauer 群中的负元，故 \([A]=-[A]\)。
\end{proof}

\begin{theorem}[Albert]\label{b2:thm:B-albert-involution}
设 \(K\) 为特征不为二的域，\(A\) 为有限维中心单 \(K\) 代数。则 \(A\) 存在第一类对合，当且仅当 \(2[A]=0\) 于 \(\operatorname{Br}(K)\)。
\end{theorem}
\begin{proof}
必要性由上一个命题得到。证明充分性时，先写 \(A\cong M_r(D)\)，其中 \(D\) 为中心除代数，令 \(d=\dim_KD\)。\(2[A]=0\) 使 \(D\otimes_KD\) 分裂，故它同构于 \(M_d(K)\)。若 \(D=K\)，转置已经给出所需对合。以下设 \(\deg D=n>1\)，因而 \(d=n^2\)。

构造一个能交换两张量因子的元素。\cref{b2:thm:csa-enveloping-isomorphism}的同构 \(D\otimes_KD^{\mathrm{op}}\cong\operatorname{End}_K(D)\) 遗忘乘法后，给出向量空间同构
\[
\operatorname{Sand}:D\otimes_KD\longrightarrow\operatorname{End}_K(D),
\qquad \operatorname{Sand}(a\otimes b)(x)=axb.
\]
令 \(g\) 为线性映射 \(x\mapsto\operatorname{Trd}_D(x)1\) 在此同构下的原像。约化迹与标量扩张相容；在分裂域上，\(D\) 变成 \(M_n\)，直接用矩阵单位计算得
\[
g=\sum_{i,j}E_{ij}\otimes E_{ji},\qquad
g^2=1,\qquad g(a\otimes b)=(b\otimes a)g.
\]
这些等式在分裂域成立，故由标量扩张的单射性在 \(K\) 上成立。上述元素称为 Goldman 元。\footnote{戈德曼（Oscar Goldman，1925-1986），美国数学家，与奥斯兰德共同将 Brauer 群的理论推广到交换环。}它在分裂后的自然张量空间上交换两因子，固定非零向量 \(v\otimes v\)，所以 \(1-g\) 不可逆。

在 \(B=D\otimes_KD\cong M_d(K)\) 中，真右理想 \((1-g)B\) 包含于一个极大右理想 \(I\)。矩阵代数的简单右模维数为 \(d\)，故 \(\dim_KI=d^2-d\)。而 \(1\otimes D\) 的每个非零元素在 \(B\) 中可逆，不能属于真右理想。因此
\[
I\cap(1\otimes D)=0,\qquad B=I\oplus(1\otimes D).
\]
对每个 \(a\in D\)，定义唯一的 \(\sigma(a)\in D\)，使 \(a\otimes1-1\otimes\sigma(a)\in I\)。投影的线性给出 \(\sigma\) 的 \(K\) 线性，且 \(\sigma(1)=1\)。因为 \(I\) 为右理想，
\[
\begin{aligned}
ab\otimes1-1\otimes\sigma(b)\sigma(a)
={}&(a\otimes1-1\otimes\sigma(a))(b\otimes1)\\
&+(b\otimes1-1\otimes\sigma(b))(1\otimes\sigma(a))\in I.
\end{aligned}
\]
唯一性遂给出 \(\sigma(ab)=\sigma(b)\sigma(a)\)。它的核为真双边理想，除代数的单性使核为零；有限维数进一步使它双射。

最后核对平方。将 \(a\otimes1-1\otimes\sigma(a)\in I\) 右乘 \(g\)，使用 Goldman 恒等式，得到 \(g(1\otimes a-\sigma(a)\otimes1)\in I\)。又因 \(1-g\in I\)，任意 \(z\in B\) 都有 \(gz-z\in I\)，所以 \(1\otimes a-\sigma(a)\otimes1\in I\)。与定义应用于 \(\sigma(a)\) 的等式相加，得 \(1\otimes(a-\sigma^2(a))\in I\cap(1\otimes D)=0\)。故 \(\sigma^2=1\)，它是 \(D\) 的第一类对合。在 \(M_r(D)\) 上逐项施加 \(\sigma\) 后转置，即得到 \(A\) 上的第一类对合。
\end{proof}

Albert 判据\footnote{阿尔伯特（Abraham Adrian Albert，1905-1972），美国数学家，研究结合代数与非结合代数，著有《结合代数》。}的这个证明采用 Goldman 元和除代数代表，见\cite[Theorem (3.1); Proposition (3.6); Theorem (3.8); Remark (3.11), pp. 31--36]{KnusMerkurjevRostTignol1998Involutions}。判据限制的是 Brauer 类的周期，并非说除代数的次数至多为二；\cref{b2:exm:period-two-index-four}已有周期二、指数四的例子。

\subsection{矩阵代数上的伴随对合}
\label{b2:subsec:B03:02}

设 \(b\) 是 \(K^n\) 上矩阵为 \(B\) 的非退化对称型或交替型。条件
\[
b(Xv,w)=b\bigl(v,\sigma_B(X)w\bigr)
\]
唯一确定 \(\sigma_B(X)=B^{-1}X^{\mathsf T}B\)，称为 \(b\) 的\term[xing-bansui-duihe]{伴随对合}[adjoint involution]。若 \(B^{\mathsf T}=\varepsilon B\)、\(\varepsilon=1\) 或 \(-1\)，直接代入可见 \(\sigma_B^2=1\)。伴随对合由整个型确定，不是把每个矩阵替换成其代数余子式组成的伴随矩阵。

\begin{theorem}\label{b2:thm:B-split-involutions}
设 \(K\) 为特征不为二的域，\(n\ge1\)。\(M_n(K)\) 上的每个第一类对合都可写为 \(\sigma_B(X)=B^{-1}X^{\mathsf T}B\)，其中 \(B\in\operatorname{GL}_n(K)\) 满足 \(B^{\mathsf T}=B\) 或 \(B^{\mathsf T}=-B\)。给定对合时，\(B\) 唯一到乘以非零标量。后一种情形只可能在 \(n\) 为偶数时出现。
\end{theorem}
\begin{proof}
将给定对合与转置复合，得到 \(M_n(K)\) 的 \(K\) 代数自同构。\cref{b2:thm:skolem-noether}说明这种自同构为内自同构，故存在可逆 \(B\) 使 \(\sigma(X)=B^{-1}X^{\mathsf T}B\)。计算平方得
\[
\sigma^2(X)=(B^{-1}B^{\mathsf T})X(B^{-1}B^{\mathsf T})^{-1}.
\]
它恒等，当且仅当 \(B^{-1}B^{\mathsf T}\) 属于矩阵代数的中心，即 \(B^{\mathsf T}=\lambda B\)。再转置一次得 \(\lambda^2=1\)，所以 \(\lambda=\pm1\)。两个矩阵给出相同对合，当且仅当它们的比与全部矩阵交换，也就相差标量。若 \(B^{\mathsf T}=-B\)，因特征不为二，它对应交替型；\cref{b2:thm:alternating-form-normal}说明非退化交替型的维数必为偶数。
\end{proof}

对称型产生的对合称为\term[zhengjiao-duihe]{正交对合}[orthogonal involution]，交替型产生的称为\term[xin-duihe]{辛对合}[symplectic involution]。这些名称也用于未分裂的中心单代数：由\cref{b2:cor:csa-algebraic-closure-splits}，在代数闭包 \(\Omega\) 上有 \(A\otimes_K\Omega\cong M_n(\Omega)\)；第一类对合延伸为 \(\sigma\otimes\operatorname{id}_\Omega\)，再看它属于哪种类型。这个定义不依赖分裂同构，因为由内自同构改变分裂同构只会改变型的基。

\begin{proposition}\label{b2:prop:B-involution-type-dimension}
设 \(K\) 为特征不为二的域，\(A\) 为有限维中心单 \(K\) 代数，\(\deg A=n\)，且 \(\sigma\) 是 \(A\) 上的第一类对合。不动 \(K\) 线性子空间
\[
\operatorname{Sym}(A,\sigma)=\bigl\{\,a\in A\mid\sigma(a)=a\,\bigr\}
\]
在正交型时的维数为 \(n(n+1)/2\)，在辛型时为 \(n(n-1)/2\)。
\end{proposition}
\begin{proof}
在分裂情形，\(X\mapsto BX\) 把不动条件分别化为 \(BX\) 对称或反对称，数独立的矩阵条目便得维数。一般情形在代数闭包上计算；线性方程 \(\sigma(a)=a\) 的系数矩阵扩域后秩不变，所以解空间维数不变。
\end{proof}

这里的 \(\operatorname{Sym}\) 表示对合的不动子空间，所取对象与对称群不同。
\symidx{\(\operatorname{Sym}(A,\sigma)\)：对合的不动子空间}

\begin{proposition}\label{b2:prop:B-unitary-adjoint}
设 \(K/k\) 是特征不为二的二次域扩张，非恒等 \(k\) 自同构记为 \(a\mapsto\bar a\)，且 \(n\ge1\)。\(M_n(K)\) 上限制到中心为这个自同构的对合，都可写为
\[
\sigma_H(X)=H^{-1}\overline X^{\mathsf T}H,
\qquad \overline H^{\mathsf T}=H,
\]
其中 \(H\) 可逆。这是 Hermitian 型 \(h(v,w)=\overline v^{\mathsf T}Hw\) 的伴随对合。
\end{proposition}
\begin{proof}
令 \(X^*=\overline X^{\mathsf T}\)。给定对合与 \(*\) 的复合在中心上为恒等，因此是 \(K\) 代数自同构。再由\cref{b2:thm:skolem-noether}，可写 \(\sigma(X)=B^{-1}X^*B\)。平方为恒等迫使 \(B^*=\lambda B\)，而再取 \(*\) 得 \(\lambda\bar\lambda=1\)。

需要找 \(a\ne0\) 满足 \(a/\bar a=\lambda\)。若 \(\lambda\ne-1\)，取 \(a=1+\lambda\)；由 \(\bar\lambda=\lambda^{-1}\) 立即得到所需等式。若 \(\lambda=-1\)，取\cref{b2:prop:B-involution-center}证明中的非零 \(u\)，满足 \(\bar u=-u\)，并令 \(a=u\)。于是 \(H=aB\) 满足 \(H^*=\bar a\lambda B=aB=H\)，而乘以标量不改变 \(\sigma_B\)。最后直接代入，得到 \(h(Xv,w)=h\bigl(v,\sigma_H(X)w\bigr)\)。
\end{proof}

第二类对合也称为\term[you-duihe]{酉对合}[unitary involution]。它没有要求 Hermitian 型正定；例如 \(H=\operatorname{diag}(1,-1)\) 也定义一个对合。正定性、范数及完备性属于\appref{b2:app:D}讨论算子代数时另加的结构。

\begin{proposition}\label{b2:prop:B-similar-forms-adjoint}
设 \(K\) 为特征不为二的域，配有一个允许恒等的对合 \(a\mapsto\bar a\)，固定域为 \(F\)。记 \(X^*=\overline X^{\mathsf T}\)，固定 \(\varepsilon\in\{1,-1\}\)，并设 \(H,H'\in\operatorname{GL}_n(K)\) 满足 \(H^*=\varepsilon H\)、\(H'^*=\varepsilon H'\)。则伴随对合 \(\sigma_H(X)=H^{-1}X^*H\) 与 \(\sigma_{H'}\) 相同，当且仅当 \(H'=cH\) 对某个 \(c\in F^\times\) 成立。更一般地，若
\[
H'=cP^*HP,\qquad c\in F^\times,\quad P\in\operatorname{GL}_n(K),
\]
则 \(\sigma_{H'}=\operatorname{Int}(P^{-1})\sigma_H\operatorname{Int}(P)\)，其中 \(\operatorname{Int}(P)(X)=PXP^{-1}\)。
\end{proposition}
\begin{proof}
乘以中心标量显然不改变伴随公式。若两个伴随对合相同，对任意 \(X\) 都有
\[
X^*H'H^{-1}=H'H^{-1}X^*.
\]
由于 \(X^*\) 取遍全部矩阵，\(H'H^{-1}\) 属于矩阵环中心，所以 \(H'=cH\)。再取共轭转置，\(H'^*=\bar c\,\varepsilon H\)，与 \(H'^*=\varepsilon H'=\varepsilon cH\) 比较得到 \(\bar c=c\)。最后直接计算
\[
\begin{aligned}
\sigma_{H'}(X)
&=P^{-1}H^{-1}(P^*)^{-1}X^*P^*HP\\
&=P^{-1}\sigma_H(PXP^{-1})P,
\end{aligned}
\]
这就是所述共轭关系。
\end{proof}

因此伴随对合通常记录型的相似类，而不记录其全部等距信息。例如 \(\mathbb R^n\)、\(n\ge1\) 上的标准正定型与它的负型给出同一个转置对合，却不能等距，因为非零向量的二次取值符号相反。命题也解释了为什么研究带对合代数的自同构时会自然遇到相似变换，而不只是等距变换。

\subsection{四元数的共轭与正交对合}
\label{b2:subsec:B03:03}

令 \(Q=(a,b)_k\)，其中 \(a,b\ne0\)，并使用 \(1,i,j,ij\) 这组基。\eqref{b2:eq:quaternion-conjugation}给出的共轭 \(\gamma\) 是第一类对合，其不动子空间只有 \(k\)。\(Q\) 的次数为二，故\cref{b2:prop:B-involution-type-dimension}用这个一维不动子空间说明 \(\gamma\) 是辛对合，即使 \(Q\) 是除代数也一样。

\begin{proposition}\label{b2:prop:B-quaternion-orthogonal-involution}
设 \(k\) 为特征不为二的域，\(a,b\in k^\times\)，\(Q=(a,b)_k\) 为以 \(i^2=a\)、\(j^2=b\)、\(ij=-ji\) 定义的四元数代数，\(\gamma\) 为其标准共轭，则映射 \(\sigma(x)=i\cdot\gamma(x)i^{-1}\) 是 \(Q\) 的正交对合，且
\[
\sigma(i)=-i,\qquad \sigma(j)=j,\qquad \sigma(ij)=ij.
\]
因此同一个四元数代数可以同时带有辛对合和正交对合。
\end{proposition}
\begin{proof}
内自同构与反同构复合仍反转乘法。由 \(ij=-ji\)，共轭作用 \(x\mapsto ixi^{-1}\) 固定 \(i\)，将 \(j,ij\) 变号；再与 \(\gamma\) 复合便得所列公式。这些公式说明 \(\sigma^2=1\)，且不动子空间为 \(k\oplus kj\oplus k(ij)\)，维数为三。\cref{b2:prop:B-involution-type-dimension}在次数二时把它识别为正交对合。
\end{proof}

共轭所对应的范数 \(x\cdot\gamma(x)\) 总在中心，见\eqref{b2:eq:quaternion-norm}；更换对合后，这个性质不能照搬。例如 \(x=1+j\) 时，\(x\cdot\sigma(x)=(1+j)^2=1+b+2j\)，一般不在 \(k\) 中。对合的类型与所使用的范数公式必须一起核对。

\begin{example}\label{b2:exm:B-quaternion-hermitian-right-module}
在这个例子中明确采用右模：令 \(D=\mathbb H\)，\(M=D^2\) 为列向量右 \(D\) 模，标量逐坐标从右相乘。定义
\[
h(x,y)=\bar x_1y_1-\bar x_2y_2
=\overline x^{\mathsf T}Jy,\qquad J=\operatorname{diag}(1,-1).
\]
这里的共轭是四元数标准共轭。它满足
\[
h(xa,yb)=\bar a\cdot h(x,y)b,\qquad
h(y,x)=\overline{h(x,y)}.
\]
若 \(h(x,y)=0\) 对全部 \(y\) 成立，依次取两个标准基向量便得 \(x=0\)，故它非退化。这是除环上的一个 Hermitian 型；公式中的标量次序是定义的一部分。
\end{example}

右 \(D\) 线性自同态由矩阵从左相乘给出，所以 \(\operatorname{End}_D(M)=M_2(D)\)。令 \(T^*=\overline T^{\mathsf T}\)，相乘仍有 \((ST)^*=T^*S^*\)：四元数共轭与转置同时反转相应的顺序。于是
\[
T^\dagger=J^{-1}T^*J,\qquad
h(Tx,y)=h(x,T^\dagger y)
\]
给出一个实线性、反转乘法且平方为恒等的运算。伴随的唯一性也可由非退化性直接得到。比如
\[
T=\begin{pmatrix}0&i\\j&0\end{pmatrix}
\quad\Longrightarrow\quad
T^\dagger=\begin{pmatrix}0&j\\i&0\end{pmatrix}.
\]
这一模型已经不同于域上的 Hermitian 型：值落在非交换除环中，不能随意交换两个标量。它说明如何在除代数上的矩阵代数中构造对合；一般带型模之间的等价理论不在这里引入。

% !TeX root = ../../Book2.tex
% 纲要标题：### B.4 自旋群的平面模型与算术现象
% 纲要位置：计划纲要.md，第 1818 行。

% 纲要要点（供目录核对）：
%
% * 自旋作用的核、实平面的倍角及有理数和有限域上的像；分开分次自同构、反序对合与 Clifford 共轭，最高外次数证明行列式为一，调用中心命题求核；正式证明反射对的归一化和五元域上全部自旋像，不对一般基域点断言满射
% <!-- 短节：不设小节 -->
%

\section{自旋群的平面模型与算术现象}
\label{b2:sec:B04}

Clifford 代数的可逆元能通过共轭作用产生正交变换。这个构造还会显露底域的影响：实数域上的满射，换成有理数点后可能不再满。本节只建立所需的代数定义和一个平面模型，不预先使用 Lie 群的覆盖空间理论。仍假设 \(\operatorname{char}k\ne2\)，且 \((V,q)\) 非退化、\(\dim V\ge1\)。

\ProseParagraphBreak
把每个生成向量固定而把乘法次序颠倒，定义 Clifford 代数上的反同构
\[
\tau(v_1\cdots v_r)=v_r\cdots v_1.
\]
其良定义可由\cref{b2:thm:clifford-universal}检验：向反代数送入同一个向量仍满足 \(v^2=q(v)1\)，故唯一延伸为反同构，且平方在生成元上为恒等。\(\tau\) 称为\term[Clifford-fanxu]{反序对合}[reversion]。它固定生成向量，但对积有 \(\tau(xy)=\tau(y)\tau(x)\)。

\ProseParagraphBreak
奇偶分次给出代数自同构 \(\alpha\)：它把生成向量变为负向量，在偶部分上为恒等，在奇部分上为负恒等，且 \(\alpha(xy)=\alpha(x)\alpha(y)\)。二者的复合 \(\alpha\circ\tau\) 称为\term[Clifford-gonge]{Clifford 共轭}[Clifford conjugation]；它也把生成向量变为负向量，却反转乘法次序。\(\alpha\) 与 \(\tau\) 在生成字上交换，所以这第三种运算仍平方为恒等，并在偶部分上与 \(\tau\) 一致。许多文献先用扭转伴随作用 \(v\mapsto\alpha(g)vg^{-1}\) 定义 Clifford 群；下述自旋群只取偶元，因而 \(\alpha(g)=g\)，扭转伴随作用与普通共轭 \(v\mapsto gvg^{-1}\) 一致。

\ProseParagraphBreak
一个偶单位的共轭要成为 \(V\) 上的变换，必须把 \(V\) 保持为自身；再要求 \(\tau(g)g=1\)，则把作用元限制在反序范数为一的元素中。这两个条件分别控制共轭的作用范围和作用元的归一化。

\begin{definition}\label{b2:def:B-spin-group}
设 \(k\) 为特征不为二的域，\((V,q)\) 为维数至少一的有限维非退化 \(k\) 二次空间，\(C=\operatorname{Cl}(V,q)\)。将 \(V\) 按\cref{b2:thm:B-clifford-basis}嵌入 \(C\)，并记 \(C^0\) 为偶部分、\(\tau\) 为固定每个 \(v\in V\) 且反转乘法次序的反序对合。若偶部分中的可逆元 \(g\) 满足 \(gVg^{-1}=V\) 和 \(\tau(g)g=1\)，就把这类元素组成的集合称为\term[zixuan-qun]{自旋群}[spin group]的 \(k\) 点集，记为
\[
\operatorname{Spin}(V,q)(k)=
\bigl\{\,g\in(C^0)^\times\ \bigm|\,
gVg^{-1}=V,\ \tau(g)g=1\,\bigr\}.
\]
这里讨论的是这个具体群；若进一步把同一构造随交换 \(k\) 代数作基变换，才进入代数群的表述。
\end{definition}
\symidx{\(\operatorname{Spin}(V,q)(k)\)：自旋群的 \(k\) 点}

\begin{proposition}\label{b2:prop:B-spin-action}
设 \(k\) 为特征不为二的域，\((V,q)\) 为维数至少一的有限维非退化 \(k\) 二次空间，\(\operatorname{Spin}(V,q)(k)\) 为上一条定义的 Clifford 代数偶部分中满足正规化及反序范数为一条件的元素集合，则该集合在乘法下成群，共轭作用给出同态
\[
\rho:\operatorname{Spin}(V,q)(k)\longrightarrow\operatorname{SO}(V,q),
\qquad \rho(g)(v)=gvg^{-1},
\]
其核为 \(\{1,-1\}\)。这里不对一般底域的 \(k\) 点断言满射。
\end{proposition}
\begin{proof}
记 \(C=\operatorname{Cl}(V,q)\)。正规化 \(V\) 的条件对乘法及逆封闭。若 \(\tau(g)g=\tau(h)h=1\)，则
\[
\tau(gh)gh=\tau(h)\tau(g)gh=1,
\qquad \tau(g^{-1})g^{-1}=1.
\]
所以所给集合是群。共轭在 \(V\) 上线性且可逆；又
\[
(gvg^{-1})^2=gv^2g^{-1}=q(v)1,
\]
故它保持 \(q\)。

取正交基 \(e_1,\ldots,e_n\)，各 \(q(e_i)\ne0\)，并令 \(\omega=e_1\cdots e_n\)。把 \(e_i\) 移过其余 \(n-1\) 个正交基向量，得到 \(e_i\omega=(-1)^{n-1}\omega e_i\)。偶单项式的每个因子都贡献这个符号，而因子个数为偶数，故总符号为一；于是所有偶元素都与 \(\omega\) 交换，特别地 \(g\omega g^{-1}=\omega\)。另一方面左侧等于 \(\rho(g)(e_1)\cdots\rho(g)(e_n)\)。Clifford 乘法的关系可能产生低次数项，不能把这种乘积本身直接当作外积；进入伴随分次以后，最高次数的类才是外积。对 \(x\in F_nC\)，记 \([x]_n\) 为它在 \(F_nC/F_{n-1}C\) 中的类经\eqref{b2:eq:B-clifford-associated-graded}对应的最高外次数分量，则
\[
\begin{aligned}
[g\omega g^{-1}]_n
&=\rho(g)(e_1)\wedge\cdots\wedge\rho(g)(e_n)\\
&=\det\rho(g)\,e_1\wedge\cdots\wedge e_n,\\
[\omega]_n&=e_1\wedge\cdots\wedge e_n.
\end{aligned}
\]
第二个等号是最高外幂上的行列式展开。由于 \(g\omega g^{-1}=\omega\) 且这个外积非零，得到 \(\det\rho(g)=1\)。

由于 \(V\) 生成 \(C\)，共轭作用在 \(V\) 上为恒等，恰好意味着作用元与整个 \(C\) 交换。因此
\[
\begin{aligned}
\ker\rho
&=\operatorname{Spin}(V,q)(k)\cap Z(C)\\
&=\{\,g\in(C^0)^\times\cap Z(C)\mid\tau(g)g=1\,\}.
\end{aligned}
\]
\cref{b2:prop:B-clifford-center}给出 \(Z(C)\cap C^0=k1\)，因此核中的元素必须是非零标量。对标量 \(g\)，反序对合保持它，归一化条件便是 \(g^2=1\)。域的特征不为二，恰有两个解 \(1,-1\)，它们都满足定义中的条件，所以核为 \(\{1,-1\}\)。
\end{proof}

作为实际计算，取 \(q=\langle1,1\rangle\)，令 \(u=e_1e_2\)，则 \(u^2=-1\)、\(\tau(u)=-u\)。所有偶元素形如 \(g=a+bu\)，且
\[
\tau(g)g=a^2+b^2.
\]
若 \(a^2+b^2=1\)，则 \(g^{-1}=a-bu\)，直接相乘得到
\begin{equation}\label{b2:eq:B-spin-plane}
\rho(a+bu)=
\begin{pmatrix}a^2-b^2&2ab\\-2ab&a^2-b^2\end{pmatrix}.
\end{equation}
因此这时正规化条件自动成立，自旋群恰由 \(a^2+b^2=1\) 的偶元素组成。在 \(k=\mathbb R\) 上令 \(a=\cos\theta\)、\(b=-\sin\theta\)，则
\[
a^2-b^2=\cos2\theta,\qquad 2ab=-\sin2\theta.
\]
所以 \(g=\cos\theta-\sin\theta\,u\) 的作用是角度 \(2\theta\) 的旋转。每个平面旋转有两个互为负数的原像，倍角现象完全来自这次 Clifford 乘法计算。

\ProseParagraphBreak
在 \(k=\mathbb Q\) 上，旋转矩阵 \(\begin{psmallmatrix}0&-1\\1&0\end{psmallmatrix}\) 保持 \(\langle1,1\rangle\)，且行列式为一，所以这个有理特殊正交变换确实存在。它却没有有理自旋原像：\eqref{b2:eq:B-spin-plane}要求 \(a^2=b^2=1/2\)，而 \(1/2\) 不是有理数的平方。所缺的是满足这些条件的有理作用元，不能把实平面的满射结论无条件搬到任意域。这与\cref{b2:exm:B-rational-norm-obstruction}的范数障碍类似：实数中可解的二次方程，在有理数中可能受到平方类或同余条件的限制。

\ProseParagraphBreak
底域也决定二次方程是否有非零解。例如 \(x^2+y^2-3z^2=0\) 在实数域有解 \((\sqrt3,0,1)\)，在有理数域却只有零解：\(z=0\) 时显然 \(x=y=0\)，而 \(z\ne0\) 时除以 \(z^2\)，就得到\cref{b2:exm:B-rational-norm-obstruction}已经排除的范数方程 \((x/z)^2+(y/z)^2=3\)。这里的模三计算提供了实数方法看不到的算术障碍。

\ProseParagraphBreak
上述偶元素的共轭还可以从反射产生：一个非各向同性向量给出一次扭转共轭，两次这样的作用复合以后，两个负号相消，便成为偶元素的普通共轭。

\begin{proposition}\label{b2:prop:B-spin-reflection-pair}
设 \(k\) 为特征不为二的域，\((V,q)\) 为维数至少一的有限维非退化 \(k\) 二次空间，\(b\) 为其归一化极化型，\(C=\operatorname{Cl}(V,q)\)。若 \(z,w\in V\) 满足 \(q(z),q(w)\ne0\)，记
\[
r_z(v)=v-\frac{2b(v,z)}{q(z)}z,
\qquad r_w(v)=v-\frac{2b(v,w)}{q(w)}w
\]
为沿这两个向量的反射，则对每个 \(v\in V\) 有
\[
-zvz^{-1}=r_z(v),\qquad
(zw)v(zw)^{-1}=r_z\bigl(r_w(v)\bigr).
\]
若进一步有 \(q(z)q(w)=c^2\)、\(c\in k^\times\)，则 \(c^{-1}zw\in\operatorname{Spin}(V,q)(k)\)，且其作用为 \(r_z\circ r_w\)。
\end{proposition}
\begin{proof}
关系 \(z^2=q(z)1\) 说明 \(z^{-1}=z/q(z)\)。利用极化关系 \(zv+vz=2b(v,z)1\)，有
\[
-zvz^{-1}
=\bigl(vz-2b(v,z)1\bigr)z^{-1}
=v-\frac{2b(v,z)}{q(z)}z.
\]
\cref{b2:lem:orthogonal-reflection}说明这个公式确是反射；对 \(w\) 同样成立。因 \((zw)^{-1}=w^{-1}z^{-1}\)，将两式复合时两个负号相消，得到
\[
r_z\bigl(r_w(v)\bigr)=z w v w^{-1}z^{-1}.
\]
所以偶单位 \(zw\) 的共轭保持 \(V\)。反序对合固定 \(z,w\)，故
\[
\tau(zw)zw=wzzw=q(z)q(w)1.
\]
当这个标量为 \(c^2\) 时，\(g=c^{-1}zw\) 为偶单位，它与 \(zw\) 有相同的共轭作用，且 \(\tau(g)g=c^{-2}q(z)q(w)1=1\)。这验证了自旋群定义的两个条件。
\end{proof}

\begin{proposition}\label{b2:prop:B-finite-spin-image}
在 \(k=\mathbb F_5\) 上，令 \(V=k^2\)、\(q(xe_1+ye_2)=x^2+y^2\)，并在 \(\operatorname{Cl}(V,q)\) 中记 \(u=e_1e_2\)。则
\[
\operatorname{Spin}(V,q)(k)=\{1,-1,u,-u\},
\qquad
\rho\bigl(\operatorname{Spin}(V,q)(k)\bigr)=\{I_2,-I_2\}.
\]
相应的特殊正交群为
\[
\operatorname{SO}(V,q)
=\left\{I_2,-I_2,
\begin{pmatrix}0&-1\\1&0\end{pmatrix},
\begin{pmatrix}0&1\\-1&0\end{pmatrix}\right\}.
\]
特别地，矩阵 \(\begin{psmallmatrix}0&-1\\1&0\end{psmallmatrix}\) 不在自旋作用的像中。
\end{proposition}
\begin{proof}
由平面模型，偶元素为 \(a+bu\)，属于自旋群恰好要求 \(a^2+b^2=1\)。\(\mathbb F_5\) 中的平方是 \(0,1,4\)，其中只有 \(1+0\) 与 \(0+1\) 等于一。因此全部解为 \((a,b)=(\pm1,0),(0,\pm1)\)，给出所列四个自旋元。\eqref{b2:eq:B-spin-plane}把 \(\pm1\) 送到 \(I_2\)，把 \(\pm u\) 送到 \(-I_2\)。

若 \(A=\begin{psmallmatrix}x&y\\z&w\end{psmallmatrix}\) 保持 \(q\) 且 \(\det A=1\)，则 \(A^{\mathsf T}A=I_2\)。将 \(A^{\mathsf T}=A^{-1}\) 与行列式为一时的逆矩阵公式比较，得到 \(w=x,y=-z\)，所以
\[
A=\begin{pmatrix}x&-z\\z&x\end{pmatrix},\qquad x^2+z^2=1.
\]
同一平方计算只给出 \((x,z)=(\pm1,0),(0,\pm1)\)，正是所列四个矩阵；每个矩阵也直接满足正交及行列式为一的条件。两个非对角矩阵不等于 \(\pm I_2\)，故不在自旋作用的像中。
\end{proof}

继续研究这些问题时，自旋群的一般结构需要代数群与群作用的知识；数域二次型则要比较实、复及非 Archimedes 完备化上的解，并研究这些信息能否决定原域上的解。范数方程、平方类、Clifford 乘法及\cref{b2:prop:B-binary-norm-classification}所给的二元分类判据，已经能够处理上面的例子。关于实二次空间上的自旋群，可进一步参阅\cite[Chapter 16]{Porteous1995CliffordClassicalGroups}。


\begin{chaptersummary}
Clifford 代数的泛性质把保持二次关系的线性映射延伸为代数同态，而有序基定理保证这些关系没有消去原来的向量空间。它按字长的伴随分次代数是外代数；在低维中，偶部分又给出二次代数与四元数。这几种描述各自保留不同的信息，不能把分次同构与原代数同构混淆。

\ProseParagraphBreak
Witt 环在忽略双曲分量以后，仍保留张量积及各向异性部分。维数奇偶和判别式是容易计算的不变量，二元型的进一步区别则由范数群决定。低维偶 Clifford 代数可以补充判别式的信息，但也未必单独决定型。对合从另一个方向记录型：矩阵转置的各种共轭形式对应双线性型，第二类对合必须保留中心上的非平凡共轭；同一个四元数代数也能配备不同类型的对合。将 Gram 矩阵乘以非零中央标量不改变伴随运算；一般的型相似还包含换基，相应伴随对合只保证共轭。因此对合并不记录全部等距信息；除环上的具体公式还须明确左右模和标量次序。

\ProseParagraphBreak
自旋群通过 Clifford 代数中的共轭作用映到特殊正交群，其核为两个标量元。在实平面上，这个映射表现为倍角；在有理数点或有限域点上，求原像还受到平方与范数条件的限制。Witt 类中的双曲分量已被忽略，伴随对合中的中央尺度已被忽略，偶 Clifford 代数也不足以单独恢复等距类；自旋作用则回答某些正交变换能否由指定底域中的 Clifford 元素产生。比较这些构造时，要逐一核对所保留的信息及实际底域。
\end{chaptersummary}

\BookExercises

\begin{exercise}\label{b2:ex:B-one-dimensional-clifford}
分别求实一维型 \(\langle1\rangle\)、\(\langle-1\rangle\) 和零型的 Clifford 代数，并用代数性质证明三者互不同构。
\end{exercise}
\begin{solution}
它们依次为 \(\mathbb R[t]/(t^2-1)\cong\mathbb R\times\mathbb R\)、\(\mathbb R[t]/(t^2+1)\cong\mathbb C\) 与 \(\mathbb R[t]/(t^2)\)。第一项的同构为 \(f\mapsto\bigl(f(1),f(-1)\bigr)\)，其逆可由一次插值得到。第三项含非零幂零元 \(t\)，前两项没有，故它不同于前两项；第二项是域，第一项含非零零因子 \((1,0),(0,1)\)，所以也不相同。
\end{solution}

\begin{exercise}\label{b2:ex:B-zero-clifford-radical}
设二维空间上 \(q=0\)。写出其 Clifford 代数的乘法关系，并证明由两个基向量生成的理想 \(J\) 满足 \(J^2\ne0\)、\(J^3=0\)。求这个代数的可逆元。
\end{exercise}
\begin{solution}
有基 \(1,e,f,ef\)，关系为 \(e^2=f^2=0,fe=-ef\)。所以 \(J=ke\oplus kf\oplus k(ef)\)，\(J^2=k(ef)\ne0\)，而再乘 \(e,f\) 均为零。元素 \(a+x\)、\(x\in J\) 在 \(a\ne0\) 时可逆，逆为 \(a^{-1}-a^{-2}x+a^{-3}x^2\)，相乘后剩余的三次项为零；在 \(a=0\) 时它幂零，不可能可逆。这里与外代数的同构是原代数同构，因为二次关系本来就是齐次的。
\end{solution}

\begin{exercise}\label{b2:ex:B-clifford-volume-center}
设实三维型 \(q=\langle1,1,-1\rangle\)，记 \(C=\operatorname{Cl}(q)\)、\(\omega=e_1e_2e_3\)。求 \(C\) 的两个非平凡中央幂等元，并利用 \(C^0\cong M_2(\mathbb R)\) 证明 \(C\cong M_2(\mathbb R)\times M_2(\mathbb R)\)。说明为什么这里偶 Clifford 代数是中心单代数，而整个 Clifford 代数不是单代数。
\end{exercise}
\begin{solution}
由\cref{b2:prop:B-clifford-center}，\(\omega\) 为中心元且 \(\omega^2=(-1)^3(1\cdot1\cdot(-1))=1\)。因此
\[
p_+=\frac{1+\omega}{2},\qquad p_-=\frac{1-\omega}{2}
\]
满足 \(p_\pm^2=p_\pm\)、\(p_+p_-=0\)、\(p_++p_-=1\)。两者均非零，因为 \(1\) 与 \(\omega\) 是不同的 Clifford 基元素。

\(\omega\) 为奇可逆元，故 \(C^1=\omega C^0\)，每个 \(x\in C\) 唯一写成 \(a+\omega b\)、\(a,b\in C^0\)。中央性与 \(\omega^2=1\) 给出代数同构
\[
C\longrightarrow C^0\times C^0,\qquad
a+\omega b\longmapsto(a+b,a-b).
\]
逆映射把 \((u,v)\) 送到 \((u+v)/2+\omega(u-v)/2\)，所以它确为双射。由\cref{b2:prop:B-low-clifford}，\(C^0\cong(-1,1)_{\mathbb R}\)。因为 \(1=N_{\mathbb C/\mathbb R}(1)\)，\cref{b2:thm:quaternion-splitting}给出 \(C^0\cong M_2(\mathbb R)\)。于是整个 \(C\) 同构于两个矩阵代数的直积，其非平凡中央幂等元恰为 \(p_+,p_-\)，并给出两个非零真双边理想；所以 \(C\) 不单，而偶部分是实中心单矩阵代数。
\end{solution}

\begin{exercise}\label{b2:ex:B-split-even-clifford}
把实三维型 \(\langle1,1,-1\rangle\) 的偶 Clifford 代数具体表示为 \(M_2(\mathbb R)\)。
\end{exercise}
\begin{solution}
由\cref{b2:prop:B-low-clifford}，它为 \((-1,1)_{\mathbb R}\)。令
\[
I=\begin{pmatrix}0&-1\\1&0\end{pmatrix},\qquad
J=\begin{pmatrix}1&0\\0&-1\end{pmatrix}.
\]
有 \(I^2=-1,J^2=1,JI=-IJ\)。\(1,J\) 生成对角矩阵，\(I,IJ\) 线性无关并生成非对角矩阵，故四者为矩阵代数的基。把 \(e_1e_2\) 送到 \(I\)、\(e_1e_3\) 送到 \(J\)，便给出所求同构。
\end{solution}

\begin{exercise}\label{b2:ex:B-graded-tensor-necessary}
用实型 \(\langle1\rangle\perp\langle1\rangle\) 说明\cref{b2:prop:B-clifford-orthogonal-sum}中的分次张量积不能换成普通张量积。
\end{exercise}
\begin{solution}
每个一维 Clifford 代数均为 \(\mathbb R\times\mathbb R\)，所以普通张量积为交换代数 \(\mathbb R^4\)。而 \(\operatorname{Cl}\bigl(\langle1,1\rangle\bigr)=(1,1)_{\mathbb R}\) 中 \(e_1e_2=-e_2e_1\ne0\)，故不交换。也可把两个生成元送到 \(\operatorname{diag}(1,-1)\) 与 \(\begin{psmallmatrix}0&1\\1&0\end{psmallmatrix}\)，直接得到 \(M_2(\mathbb R)\)。两种张量积虽有相同的向量空间维数，乘法却不同。
\end{solution}

\begin{exercise}\label{b2:ex:B-witt-zero-divisors}
证明 \(W(\mathbb Q)\) 中 \(\bigl[\langle1,-2\rangle\bigr]\) 与 \(\bigl[\langle1,2\rangle\bigr]\) 都非零，但乘积为零。
\end{exercise}
\begin{solution}
\(x^2-2y^2=0\) 在有理数域只有零解，因为二不是有理数平方；\(x^2+2y^2=0\) 也只有零解。两型均为非零各向异性型，由\cref{b2:thm:witt-group}的各向异性代表唯一性，其类非零。乘积型是 \(\langle1,2,-2,-4\rangle\)。中间两项构成双曲平面，末项把基向量除以二后为 \(-1\)，与首项也构成双曲平面。因此乘积为 \(H\perp H\) 的类，即零。
\end{solution}

\begin{exercise}\label{b2:ex:B-discriminant-insufficient}
在 \(W(\mathbb R)\) 中，证明正定四维型的类属于 \(I(\mathbb R)^2\)，判别式为一，却不等于零。
\end{exercise}
\begin{solution}
正定四维型是 \(\langle1,1\rangle\otimes\langle1,1\rangle\)，两个因子的维数均为偶数，所以其类属于基本理想的平方。判别式为 \((-1)^6=1\)。它的符号差为四，由 \(W(\mathbb R)\cong\mathbb Z\) 对应非零整数四，故不是零类。这也给出判别式为一不意味着双曲的具体反例。
\end{solution}

\begin{exercise}\label{b2:ex:B-binary-norm-isometries}
在 \(\mathbb Q\) 上分别构造 \(\langle5,5\rangle\to\langle1,1\rangle\) 与 \(\langle7,-14\rangle\to\langle1,-2\rangle\) 的等距映射。
\end{exercise}
\begin{solution}
第一项利用 \(N_{\mathbb Q(i)/\mathbb Q}(1+2i)=5\)，取矩阵 \(P=\begin{psmallmatrix}1&-2\\2&1\end{psmallmatrix}\)，便有 \(P^{\mathsf T}P=5I_2\) 且 \(\det P=5\ne0\)。第二项利用 \(N_{\mathbb Q(\sqrt2)/\mathbb Q}(3+\sqrt2)=7\)，乘以 \(3+\sqrt2\) 的矩阵为 \(Q=\begin{psmallmatrix}3&2\\1&3\end{psmallmatrix}\)。直接计算
\[
Q^{\mathsf T}\begin{pmatrix}1&0\\0&-2\end{pmatrix}Q
=\begin{pmatrix}7&0\\0&-14\end{pmatrix},\qquad \det Q=7.
\]
因此两矩阵均给出指定方向的等距双射。
\end{solution}

\begin{exercise}\label{b2:ex:B-witt-extension}
说明把实二次型扩张到复数域诱导一个环同态 \(W(\mathbb R)\to W(\mathbb C)\)，并求它在整数描述下的形式与核。
\end{exercise}
\begin{solution}
扩域与正交直和、张量积相容，并把实双曲平面变成复双曲平面，故对 Witt 类良定义且是环同态。实型惯性为 \((p_+,p_-)\) 时，其实 Witt 类对应 \(p_+-p_-\)，扩域后的复类对应 \(p_++p_-\) 模二；两整数同余。因此同态是 \(\mathbb Z\to\mathbb Z/2\mathbb Z\) 的约化映射，核为 \(2\mathbb Z=I(\mathbb R)\)。
\end{solution}

\begin{exercise}\label{b2:ex:B-diagonal-adjoint}
对 \(B=\operatorname{diag}(1,d)\)、\(d\in k^\times\)，计算 \(M_2(k)\) 上的 \(\sigma_B\)，并求其不动子空间与满足 \(\sigma_B(X)=-X\) 的子空间。
\end{exercise}
\begin{solution}
有
\[
\sigma_B\begin{pmatrix}a&b\\c&e\end{pmatrix}
=\begin{pmatrix}a&dc\\b/d&e\end{pmatrix}.
\]
不动子空间由 \(E_{11},E_{22},dE_{12}+E_{21}\) 生成，维数三；反号子空间由 \(-dE_{12}+E_{21}\) 生成，维数一。由于二可逆，每个矩阵唯一分解为 \(\bigl(X+\sigma_B(X)\bigr)/2\) 与 \(\bigl(X-\sigma_B(X)\bigr)/2\) 的和，与这两个维数相符。
\end{solution}

\begin{exercise}\label{b2:ex:B-symplectic-adjugate}
此处伴随矩阵指代数余子式矩阵的转置，即 classical adjugate。取 \(J=\begin{psmallmatrix}0&1\\-1&0\end{psmallmatrix}\)。证明 \(\sigma_J(X)=\operatorname{tr}(X)I_2-X\)，它恰好是二阶矩阵的伴随矩阵。解释为什么这种等同不能直接推广到三阶矩阵。
\end{exercise}
\begin{solution}
若 \(X=\begin{psmallmatrix}a&b\\c&d\end{psmallmatrix}\)，直接相乘得到 \(J^{-1}X^{\mathsf T}J=\begin{psmallmatrix}d&-b\\-c&a\end{psmallmatrix}\)，同时等于所列两种表达式。它对 \(X\) 线性、反转乘法且平方为恒等。三阶伴随矩阵的各项是二次多项式，通常不是线性映射；例如在 \(\mathbb Q\) 上，\(\operatorname{adj}(2I_3)=4I_3\ne2\operatorname{adj}(I_3)\)。此外非退化交替型不存在于三维，不能照用二阶辛模型。
\end{solution}

\begin{exercise}\label{b2:ex:B-indefinite-unitary}
在 \(M_2(\mathbb C)\) 上取 \(H=\operatorname{diag}(1,-1)\)。求 \(\sigma_H\) 的不动矩阵，找出一个满足 \(\sigma_H(X)=X\)、\(X^2=-I_2\) 的矩阵。说明这不与正定内积的自伴算子性质矛盾。
\end{exercise}
\begin{solution}
计算给出 \(\sigma_H\left(\begin{psmallmatrix}a&b\\c&d\end{psmallmatrix}\right)=\begin{psmallmatrix}\bar a&-\bar c\\-\bar b&\bar d\end{psmallmatrix}\)。不动条件为 \(a,d\in\mathbb R\)、\(c=-\bar b\)。取 \(X=\begin{psmallmatrix}0&1\\-1&0\end{psmallmatrix}\) 即满足要求。这里的 Hermitian 型不正定，\(h(Xv,Xv)=h(v,X^2v)=-h(v,v)\) 完全可能；正定内积下不能发生的现象，不妨碍它成为合法的酉对合。
\end{solution}

\begin{exercise}\label{b2:ex:B-no-odd-division-involution}
设 \(k\) 为特征不为二的域，\(D\ne k\) 为奇素数次数 \(\ell\) 的中心除 \(k\) 代数。若有限扩张 \(E/k\) 满足 \(\gcd([E:k],\ell)=1\)，证明 \(D\otimes_kE\) 仍为中心除代数，并证明它和 \(M_2(D\otimes_kE)\) 都没有第一类对合。后一代数的次数是偶数，为什么仍有这个障碍？
\end{exercise}
\begin{solution}
令 \(n=[E:k]\)、\(d=\operatorname{ind}(D\otimes_kE)\)。\cref{b2:prop:index-after-base-change}给出
\[
d\mid\ell,\qquad \ell\mid nd.
\]
由于 \(n\) 与 \(\ell\) 互素，第二式迫使 \(\ell\mid d\)，所以 \(d=\ell\)。扩域后代数的次数也为 \(\ell\)，故在其 Artin--Wedderburn 分解 \(M_r(D')\) 中，有 \(\ell=r\deg D'\) 和 \(\deg D'=d=\ell\)，得到 \(r=1\)。因此它仍为除代数，且次数大于一，Brauer 类非零。

矩阵代数 \(M_2(D\otimes_kE)\) 与这个除代数有相同的非零 Brauer 类和奇数指数 \(\ell\)。\cref{b2:cor:B-odd-index-no-involution}同时排除两者的第一类对合。后一代数虽有次数 \(2\ell\)，但加矩阵块没有改变 Brauer 类；Albert 判据约束的是该类的周期，不是矩阵次数的奇偶。
\end{solution}

\begin{exercise}\label{b2:ex:B-quaternion-fixed-norm}
令 \(Q=(a,b)_k\)，并取\cref{b2:prop:B-quaternion-orthogonal-involution}中的 \(\sigma\)。计算四元数范数在 \(\operatorname{Sym}(Q,\sigma)\) 上的限制，并与 \(\langle1,1,1\rangle\) 的偶 Clifford 代数联系起来。
\end{exercise}
\begin{solution}
不动元素形如 \(x+yj+zij\)，故\eqref{b2:eq:quaternion-norm}给出限制型 \(x^2-by^2+abz^2\)。在 \(Q=\mathbb H\)、\(a=b=-1\) 时，它是正定三维型。反过来，\cref{b2:prop:B-low-clifford}将这个型的偶 Clifford 代数识别为 \((-1,-1)_{\mathbb R}=\mathbb H\)。这里先在指定正交对合的不动子空间上限制四元数范数，再使用三维偶 Clifford 公式恢复这个四元数代数。所算的是一个三维限制型，原四维范数依然保留；这一计算本身也没有给出一般带对合代数与二次型之间的分类对应。
\end{solution}

\begin{exercise}\label{b2:ex:B-rational-spin-parameters}
求 \(a^2+b^2=1\) 的全部有理解，并利用\eqref{b2:eq:B-spin-plane}写出相应的有理自旋作用。计算参数 \(t=1/2\) 的一个例子。
\end{exercise}
\begin{solution}
除 \((-1,0)\) 外，令 \(t=b/(1+a)\)，由圆方程解出
\[
a=\frac{1-t^2}{1+t^2},\qquad b=\frac{2t}{1+t^2},\qquad t\in\mathbb Q.
\]
反向代入也满足圆方程，所以没有遗漏。把这两个表达式代入\eqref{b2:eq:B-spin-plane}即可得全部像。取 \(t=1/2\)，有 \(a=3/5,b=4/5\)，因此
\[
\rho(3/5+4u/5)=\frac1{25}\begin{pmatrix}-7&24\\-24&-7\end{pmatrix}.
\]
矩阵列向量平方和为一，互相正交，行列式也为一；这些等式可直接在有理数中核验。
\end{solution}

\begin{exercise}\label{b2:ex:B-clifford-reflection-product}
在 \(\mathbb Q^2\) 上取 \(q=\langle1,1\rangle\)，令 \(z=e_1\)、\(w=5e_1+12e_2\)、\(u=e_1e_2\)。在原平面坐标中计算沿 \(w\)、再沿 \(z\) 的反射复合，并构造一个产生这次变换的有理自旋元素。证明它只有两个自旋原像。
\end{exercise}
\begin{solution}
\(q(z)=1\)、\(q(w)=169\)，故\cref{b2:prop:B-spin-reflection-pair}给出
\[
g=\frac1{13}zw=\frac5{13}+\frac{12}{13}u\in\operatorname{Spin}(q)(\mathbb Q).
\]
反射公式给出
\[
r_z=\begin{pmatrix}-1&0\\0&1\end{pmatrix},\qquad
r_w=I_2-\frac2{169}\begin{pmatrix}25&60\\60&144\end{pmatrix}
=\frac1{169}\begin{pmatrix}119&-120\\-120&-119\end{pmatrix}.
\]
所以指定顺序的复合为
\[
r_zr_w=\frac1{169}\begin{pmatrix}-119&120\\-120&-119\end{pmatrix}.
\]
把 \(a=5/13\)、\(b=12/13\) 代入\eqref{b2:eq:B-spin-plane}，也得到这个矩阵，因而核验了 Clifford 共轭与两次反射的次序。若另一个自旋元素 \(h\) 给出同一变换，\cref{b2:prop:B-spin-action}说明 \(h^{-1}g\in\{1,-1\}\)，故原像恰为 \(g,-g\)。
\end{solution}

\begin{exercise}\label{b2:ex:B-finite-spin-image}
在 \(k=\mathbb F_7\)、\(q=\langle1,1\rangle\) 中列出全部自旋元素及其像，计算 \(\operatorname{SO}(q)(k)\) 的阶，并证明 \(\begin{pmatrix}2&2\\-2&2\end{pmatrix}\) 属于该群却没有 \(k\) 自旋原像。
\end{exercise}
\begin{solution}
\(\mathbb F_7\) 的平方为 \(0,1,2,4\)。\(a^2+b^2=1\) 的解恰为
\[
(\pm1,0),\quad(0,\pm1),\quad(2,2),(2,-2),(-2,2),(-2,-2).
\]
因此自旋群有八个元素 \(a+bu\)，其中 \(u^2=-1\)。\eqref{b2:eq:B-spin-plane}将前两类分别送到 \(I_2,-I_2\)；后四个解满足 \(a^2-b^2=0\)、\(2ab=\pm1\)，给出两个非对角矩阵。因此全部八个自旋元素的像为
\[
\left\{I_2,-I_2,
\begin{pmatrix}0&1\\-1&0\end{pmatrix},
\begin{pmatrix}0&-1\\1&0\end{pmatrix}\right\}.
\]
对二维矩阵，正交条件与行列式一使逆矩阵同时等于转置及代数余子式的转置。逐条比较，得到全部特殊正交矩阵恰为
\[
\begin{pmatrix}s&t\\-t&s\end{pmatrix},\qquad s^2+t^2=1.
\]
已列出的同一圆方程共有八个解，故特殊正交群的阶也是八。\(s=t=2\) 满足 \(s^2+t^2=8=1\) 于 \(\mathbb F_7\)，所以题给矩阵在群中；但它不属于上述四元像。这也可从倍角方程看出：若 \(a^2+b^2=1\)、\(a^2-b^2=2\)，则 \(a^2=3/2=5\)，而五不是该域中的平方。
\end{solution}

\begin{exercise}\label{b2:ex:B-clifford-distinguishes-ternary}
设 \(k\) 为特征不为二的域，\(q=\langle a,b,c\rangle\)、\(a,b,c\in k^\times\)，并取 \(\lambda\in k^\times\)。证明 \(q\) 与 \(\lambda q\) 的偶 Clifford 代数总同构，而两型等距当且仅当 \(\lambda\) 为平方。用 \(k=\mathbb Q\)、\(q=\langle1,1,1\rangle\)、\(\lambda=3\) 给出一个正定型之间的例子。
\end{exercise}
\begin{solution}
\cref{b2:prop:B-low-clifford}分别给出
\[
\operatorname{Cl}^0(q)\cong(-ab,-ac)_k,\qquad
\operatorname{Cl}^0(\lambda q)\cong(-\lambda^2ab,-\lambda^2ac)_k.
\]
把后一四元数的两个生成元各除以 \(\lambda\)，就得到前一代数的平方与反交换关系；四个基元素仍线性无关，因此这是代数同构。

三维 Gram 矩阵整体乘以 \(\lambda\)，使行列式乘以 \(\lambda^3\)，而等距换基只能使行列式乘以一个平方。所以等距迫使 \(\lambda^3\)、从而 \(\lambda\) 为平方。反过来，若 \(\lambda=s^2\)，乘 \(s\) 就给出从 \((k^3,\lambda q)\) 到 \((k^3,q)\) 的等距双射。

在所给有理例子中，两型都为正定，偶 Clifford 代数都同构于 \((-1,-1)_{\mathbb Q}\)，但三不是有理数平方，故不等距；其判别式平方类分别为 \(-1\) 与 \(-3\)。这里即使实符号相同，偶 Clifford 同构仍不能消除有理等距障碍。
\end{solution}

\begin{exercise}\label{b2:ex:B-similarity-not-isometry}
在 \(\mathbb Q^2\) 上取 \(q=\langle1,1\rangle\)、\(r=\langle13,13\rangle\)、\(s=\langle7,7\rangle\)。比较三个型的伴随对合与等距子群，构造 \(r\to q\) 的等距，并用模七计算证明 \(s\) 与 \(q\) 不等距。
\end{exercise}
\begin{solution}
三个 Gram 矩阵都是 \(I_2\) 的非零有理标量倍，所以都给出转置对合；条件 \(P^{\mathsf T}(aI_2)P=aI_2\) 也都等价于 \(P^{\mathsf T}P=I_2\)，因此等距子群相同。取
\[
P=\begin{pmatrix}2&-3\\3&2\end{pmatrix},\qquad
P^{\mathsf T}P=13I_2,\qquad\det P=13,
\]
便得到所要求的等距 \(r\to q\)，它来自范数 \(N_{\mathbb Q(i)/\mathbb Q}(2+3i)=13\)。

若 \(s\cong q\)，则 \(q\) 表示七，即有 \(x^2+y^2=7\) 的有理解。清分母并约去公因子后得到 \(X^2+Y^2=7Z^2\)，其中 \(X,Y,Z\) 为互素整数、\(Z\ne0\)。模七的非零平方为 \(1,2,4\)，两个平方之和为零只能是二者均为零，故 \(7\mid X,Y\)。写 \(X=7X_1\)、\(Y=7Y_1\)，代回得 \(Z^2=7(X_1^2+Y_1^2)\)，于是 \(7\mid Z\)，与互素矛盾。因此三个相同的伴随对合和相同的等距子群，并没有强迫三个型互相等距；具体等距仍须判断相应范数。
\end{solution}

\begin{exercise}\label{b2:ex:B-quaternion-hermitian-isometries}
在右 \(\mathbb H\) 模 \(\mathbb H^2\) 上取标准正定 Hermitian 型 \(h(x,y)=\bar x_1y_1+\bar x_2y_2\)。对 \(a,b\in\mathbb H^\times\)，令 \(T=\operatorname{diag}(a,b)\) 从左作用。求 \(T^\dagger\)，并给出 \(T\) 为等距变换或相似变换的充要条件。
\end{exercise}
\begin{solution}
Gram 矩阵为单位，故 \(T^\dagger=T^*=\operatorname{diag}(\bar a,\bar b)\)，从而 \(T^\dagger T=\operatorname{diag}(N(a),N(b))\)。逐坐标代入配对，得到等距当且仅当 \(N(a)=N(b)=1\)。相似条件为 \(h(Tx,Ty)=c\cdot h(x,y)\)、\(c\in\mathbb R^\times\)，当且仅当 \(N(a)=N(b)=c\)。非零实四元数的范数为正，故这个乘子必为正。例如 \(a=i,b=j\) 给出等距变换，\(a=1+i,b=1+j\) 则给出乘子为二的相似变换。
\end{solution}
